% 2.1.2
% Ficheiro incluído por main.tex

\subsection[Ubuntu: IP estático e DNS do Windows]{Ubuntu: IP estático e DNS do Windows}

\esquema{2-1-2-rede-dns}{Rede e DNS do Ubuntu apontados para o Windows Server}

\textbf{Objetivo:} no fim, o Ubuntu tem IP estático, um único perfil de rede, usa \textbf{apenas} o DC como DNS e resolve \texttt{isec.local} pelo nome. Sem o DNS apontado ao Windows, o Kerberos e o SSSD não conseguem descobrir o domínio.

\passo{Passo 1 - Iniciar a VM e abrir o terminal}

Trabalha \textbf{na consola da VM} (não por SSH): a rede vai ser reconfigurada e uma sessão SSH cairia.

\begin{simplesbox}
Confirma o nome da placa com \texttt{ip -br a}. Se não for \texttt{ens18}, troca \texttt{ens18} pelo nome real nos comandos desta secção.
\end{simplesbox}

\passo{Passo 2 - Ver a configuração de rede existente}

\emph{Objetivo:} perceber que ficheiros e ligações já configuram a placa, para evitar conflitos.

\begin{lstlisting}[style=bashstyle]
ls /etc/netplan
nmcli -f NAME,DEVICE,TYPE connection show
\end{lstlisting}

\begin{itemize}
\item \texttt{ls /\allowbreak{}etc/\allowbreak{}netplan} lista os ficheiros do Netplan (o sistema que define a rede no Ubuntu).
\item \texttt{nmcli ... connection show} lista as ligações geridas pelo NetworkManager.
\end{itemize}

\emph{Porquê:} o erro \texttt{Conflicting default route} aparece quando mais de um ficheiro ou perfil configura a mesma placa com uma rota por defeito.

\passo{Passo 3 - Criar uma pasta de backup}

\emph{Objetivo:} guardar a configuração antiga em vez de a apagar, para poder repor se algo correr mal.

\begin{lstlisting}[style=bashstyle]
sudo mkdir -p /root/netplan-bak
\end{lstlisting}

\texttt{-p} não dá erro se a pasta já existir; fica em \texttt{/\allowbreak{}root}, só acessível ao administrador.

\emph{Se saltares este passo, o Passo 4 falha com \texttt{Ficheiro ou pasta inexistente}.}

\passo{Passo 4 - Remover os ficheiros em conflito}

\emph{Objetivo:} deixar um único método de configuração de rede: o NetworkManager por linha de comandos.

\begin{lstlisting}[style=bashstyle]
cd /etc/netplan
sudo mv 00-installer-config.yaml 01-netcfg.yaml 90-NM-*.yaml /root/netplan-bak/
ls /etc/netplan
cd ~
\end{lstlisting}

\emph{Verificação esperada:} o \texttt{ls} mostra \textbf{só} \texttt{01-network-manager-all.yaml} (que apenas diz que o NetworkManager gere a rede e não configura nenhuma placa).

\emph{Nota:} os nomes dos ficheiros podem variar - a regra é mover \textbf{todos exceto} \texttt{01-network-manager-all.yaml}.

\passo{Passo 5 - Aplicar a limpeza}

\begin{lstlisting}[style=bashstyle]
sudo netplan apply
\end{lstlisting}

Como já não há ficheiro a configurar a \texttt{ens18}, o perfil antigo desaparece.

\emph{Nota:} podem aparecer avisos amarelos sobre o \texttt{systemd-networkd} - são inofensivos, porque não o usamos.

\passo{Passo 6 - Garantir que o NetworkManager está a correr}

\begin{lstlisting}[style=bashstyle]
sudo systemctl start NetworkManager
systemctl is-active NetworkManager
\end{lstlisting}

\emph{Verificação esperada:} \texttt{active}.

\emph{Porquê:} sem o serviço ativo, o \texttt{nmcli} falha com \texttt{O NetworkManager não está em execução} e a ligação do Passo 7 nunca é criada.

\passo{Passo 7 - Criar a ligação com IP estático e DNS do Windows Server}

\begin{lstlisting}[style=bashstyle]
sudo nmcli connection add type ethernet ifname ens18 con-name labX \
     ipv4.method manual \
     ipv4.addresses <IP_CLIENTE>/24 \
     ipv4.gateway 192.168.X.1 \
     ipv4.dns <IP_DC> \
     ipv4.dns-search isec.local \
     ipv4.ignore-auto-dns yes
\end{lstlisting}

\begin{center}\small
\begin{tabularx}{\linewidth}{|>{\raggedright\arraybackslash}X|>{\raggedright\arraybackslash}X|>{\raggedright\arraybackslash}X|}
\hline
\rowcolor{azul!10}\textbf{Parâmetro} & \textbf{Valor} & \textbf{O que faz} \\
\hline
\texttt{ifname} & \texttt{ens18} & Placa de rede a configurar \\
\hline
\texttt{con-name} & \texttt{labX} & Nome do perfil \\
\hline
\texttt{ipv4.method} & \texttt{manual} & IP fixo, sem DHCP \\
\hline
\texttt{ipv4.addresses} & \texttt{\textless{}IP\_CLIENTE\textgreater{}/\allowbreak{}24} & IP do Ubuntu e máscara \\
\hline
\texttt{ipv4.gateway} & \texttt{192.168.X.1} & Saída da rede \\
\hline
\texttt{ipv4.dns} & \texttt{\textless{}IP\_DC\textgreater{}} & DNS = Windows Server \\
\hline
\texttt{ipv4.dns-search} & \texttt{isec.local} & Domínio de pesquisa (escreves \texttt{SR2-SRV} em vez de \texttt{SR2-SRV.isec.local}) \\
\hline
\texttt{ipv4.ignore-auto-dns} & \texttt{yes} & Impede que outro DNS entre automaticamente \\
\hline
\end{tabularx}
\end{center}

\begin{importantebox}[title={Importante}]
\textbf{Só o DNS do Windows:} não ponhas \texttt{8.8.8.8} como alternativo. O Ubuntu pode perguntar-lhe e não encontrar os registos internos do AD. A internet continua a funcionar através do \textbf{forwarder} do DNS do Windows (1.3, Passo 7).
\end{importantebox}

\emph{Verificação esperada:} o comando termina sem erros.

\passo{Passo 8 - Ativar a ligação e confirmar}

\begin{lstlisting}[style=bashstyle]
sudo nmcli connection up labX
nmcli -f NAME,DEVICE,TYPE connection show
\end{lstlisting}

\emph{Verificação esperada:} o perfil \texttt{labX} aparece associado à placa \texttt{ens18}.

\emph{Nota:} se aparecer outro perfil antigo sem placa (ex.: \texttt{Ligação com fios 1}), apaga-o para não confundir:

\begin{lstlisting}[style=bashstyle]
sudo nmcli connection delete 'NOME'
\end{lstlisting}

\passo{Passo 9 - Confirmar o IP e a rota}

\begin{lstlisting}[style=bashstyle]
ip -br a
ip route
\end{lstlisting}

\emph{Verificação esperada:}

\begin{itemize}
\item \texttt{ens18} aparece \textbf{UP} com \texttt{\textless{}IP\_CLIENTE\textgreater{}/\allowbreak{}24} - e \textbf{não} com um endereço \texttt{169.254.x.x} (que significa que o IP estático não foi aplicado);
\item existe \textbf{uma única} linha \texttt{default via 192.168.X.1} (o conflito do início desapareceu).
\end{itemize}

\passo{Passo 10 - Confirmar que o DNS é o do Windows}

\begin{lstlisting}[style=bashstyle]
resolvectl status ens18
\end{lstlisting}

\emph{Verificação esperada:} \texttt{Current DNS Server: \textless{}IP\_DC\textgreater{}}, \texttt{DNS Servers: \textless{}IP\_DC\textgreater{}} (sem \texttt{8.8.8.8} nem \texttt{1.1.1.1}), \texttt{DNS Domain: isec.local}.

\emph{Porquê:} se ainda aparecerem DNS públicos, o Ubuntu pode perguntar a um deles e falhar a descoberta do domínio. Volta ao Passo 7 e confirma \texttt{ipv4.ignore-auto-dns yes}.

\passo{Passo 11 - Testar a ligação ao servidor}

\begin{lstlisting}[style=bashstyle]
ping -c 3 <IP_DC>
\end{lstlisting}

\emph{Verificação esperada:} \texttt{0\% packet loss}.

\emph{Porquê:} prova que há rede, mas ainda não que o DNS funciona (o ping por IP não precisa de DNS).

\passo{Passo 12 - Testar o DNS do domínio}

\emph{Objetivo:} provar que o DNS do Windows responde sobre o domínio e anuncia os serviços de autenticação.

\begin{lstlisting}[style=bashstyle]
dig @<IP_DC> isec.local +short
dig @<IP_DC> _ldap._tcp.dc._msdcs.isec.local SRV +short
\end{lstlisting}

\emph{Verificação esperada:}

\begin{itemize}
\item o 1.º devolve \texttt{\textless{}IP\_DC\textgreater{}};
\item o 2.º devolve algo como \texttt{0 100 389 sr2-srv.isec.local.}
\end{itemize}

\emph{Porquê:} os registos SRV são a peça mais importante - é com eles que o Kerberos e o SSSD descobrem o controlador de domínio. \textbf{Se falhar, é DNS: não avances.}

\passo{Passo 13 - Corrigir o ping a \texttt{isec.local} (desativar o mDNS para \texttt{.local})}

\emph{Objetivo:} fazer o Ubuntu perguntar ao DNS do Windows quando o nome acaba em \texttt{.local}.

\begin{lstlisting}[style=bashstyle]
sudo cp /etc/nsswitch.conf /etc/nsswitch.conf.bak
grep ^hosts /etc/nsswitch.conf
sudo sed -i 's/mdns4_minimal \[NOTFOUND=return\] //' /etc/nsswitch.conf
grep ^hosts /etc/nsswitch.conf
\end{lstlisting}

\emph{Porquê:} por defeito, o Ubuntu trata nomes \texttt{.local} como \textbf{mDNS} (descoberta na rede local) e responde logo "nome desconhecido", sem perguntar ao DNS do Windows. Por isso o \texttt{dig} funciona e o \texttt{ping isec.local} falha. \textbf{Este passo é obrigatório:} sem ele, o Kerberos, o \texttt{smbclient} e o \texttt{realm} também não conseguem resolver \texttt{SR2-SRV.isec.local}.

\emph{Verificação esperada:} a linha passa de \texttt{hosts: files mdns4\_minimal [NOTFOUND=return] dns} para \texttt{hosts: files dns}.

\passo{Passo 14 - Testar o \texttt{isec.local}}

\begin{lstlisting}[style=bashstyle]
getent hosts isec.local
ping -c 3 isec.local
\end{lstlisting}

\emph{Verificação esperada:} \texttt{getent} devolve \texttt{\textless{}IP\_DC\textgreater{} isec.local} e o ping responde com \texttt{0\% packet loss}.

\problemas{\hyperref[ec-b1-rede]{Erros Comuns - 2.1 Rede e DNS do cliente}}
